\documentclass[titlepage, UTF8,9pt,a4paper]{ctexart}
\usepackage{amsmath}
\usepackage{mathtools}
\DeclarePairedDelimiter\ceil{\lceil}{\rceil}
\DeclarePairedDelimiter\floor{\lfloor}{\rfloor}

\usepackage{listings}
\renewcommand{\lstlistingname}{代码块}% Listing -> 代码块

\usepackage{fontspec}
%\setmainfont{SimSun}
\usepackage{lipsum}
\usepackage{graphicx}
\usepackage{color}
%\usepackage{times}
\usepackage[colorlinks, linkcolor=black, anchorcolor=black, citecolor=black]{hyperref}
\usepackage{booktabs}
\usepackage[table, dvipsnames]{xcolor}

\usepackage{geometry} %设置页边距的宏包
\geometry{left=2.5cm,right=2.5cm,top=2.5cm,bottom=2.5cm} %设置页边距
\usepackage{setspace}   %行间距
\setstretch{1.25} % 任意行距

\usepackage{fancyhdr}%导入fancyhdr包(页眉页脚)
\usepackage{lastpage}
\fancyhead{} % 初始化页眉
\fancyhead[L]{page \thepage \  of \pageref{LastPage}}
\fancyhead[R]{\textsl{\leftmark}}
\setlength{\headheight}{14.49998pt}

\usepackage{minted}
\usepackage[fencedCode, pipeTables, headerAttributes, relativeReferences,
	    smartEllipses, theme=witiko/dot, hybrid, tightLists=false]{markdown}

\setkeys{Gin}{
  % width=\columnwidth,
  % % height=0.2\paperheight,
  keepaspectratio}
\usepackage{fancyvrb}
% redefine \VerbatimInput
\RecustomVerbatimCommand{\VerbatimInput}{VerbatimInput} {
	fontsize=\footnotesize,
	labelposition=topline,
	frame=lines,  % top and bottom rule only
	framesep=1em, % separation between frame and text
	rulecolor=\color{Gray},
}
\newcommand{\verbatiminput}[1]{%
  \VerbatimInput[
    label=\fbox{\color{Black}\texttt{\detokenize{#1}}},
  ]{#1}%
}
\global\def\markdownRendererImage#1#2#3#4{
\begin{figure}[htpb]
	\centering
	\includegraphics[width=#4\linewidth]{#3}
	\caption{#1}%
	\label{#1}
\end{figure}
}

% 设置caption字体大小
\usepackage{caption}
\captionsetup{font=footnotesize}

\title{实验标题}
\newcommand{\Supervisor}{\textsl{魏更宇副教授}}
\newcommand{\Course}{\textsl{编译原理与技术程序设计}}

% 作者
\author{
2019114514班 \quad 田所浩二 \quad 2019114514\\
2019114514班 \quad 田所浩二 \quad 2019114514\\
2019114514班 \quad 田所浩二 \quad 2019114514
}

% 日期
\date{\today}

\newcommand{\chap}[1]{\input{chap/#1}}

% 代码块格式设置
\lstset {
	frame=single, % 单线框边框
	tabsize=2, % tab显示空格数
	basicstyle=\ttfamily, % 等宽字体
	breaklines=true, % 自动换行
	flexiblecolumns, % 反正加上它
	captionpos=b, % caption 位于底部
}

\begin{document}

\pagestyle{empty}
\setlength{\parskip}{0.5em}

\begin{titlepage}
	\newcommand{\HRule}{\rule{\linewidth}{0.5mm}}
	\includegraphics[width=8cm]{title/logo.png}\\[1cm] 
	\center 
	\quad\\[1.5cm]
	\textsl{\Large 北京邮电大学 }\\[0.5cm] 
	\textsl{\large 计算机学院(国家示范性软件学院)}\\[0.5cm] 
	\makeatletter

	\HRule \\[0.4cm]
	{ \huge \bfseries \@title}\\[0cm] 
	\HRule \\[1.5cm]

	\large Author:\\[0.5cm]
	\large \@author\\[3cm]

	\vfill 
	\begin{minipage}{0.4\textwidth}
		\begin{flushleft} \large
			\emph{Course:}\\
			\Course
		\end{flushleft}
	\end{minipage}
	~
	\begin{minipage}{0.4\textwidth}
		\begin{flushright} \large
			\emph{Supervisor:} \\
			\Supervisor
		\end{flushright}
	\end{minipage}\\[1.5cm]

	\large \today\\[2cm] 
\end{titlepage}


\clearpage
\tableofcontents % 生成目录
\clearpage

\pagestyle{fancy}
\setcounter{page}{1}

\chap{chap1}
\chap{chap2}
\chap{chap3}

\end{document}

